\documentclass[otchet]{../SCWorks}
% Тип обучения (одно из значений):
%    bachelor   - бакалавриат (по умолчанию)
%    spec       - специальность
%    master     - магистратура
% Форма обучения (одно из значений):
%    och        - очное (по умолчанию)
%    zaoch      - заочное
% Тип работы (одно из значений):
%    coursework - курсовая работа (по умолчанию)
%    referat    - реферат
%  * otchet     - универсальный отчет
%  * nirjournal - журнал НИР
%  * digital    - итоговая работа для цифровой кафедры
%    diploma    - дипломная работа
%    pract      - отчет о научно-исследовательской работе
%    autoref    - автореферат выпускной работы
%    assignment - задание на выпускную квалификационную работу
%    review     - отзыв руководителя
%    critique   - рецензия на выпускную работу

% * Добавлены вручную. За вопросами к @mchernigin
\usepackage{../preamble}

\begin{document}

% Кафедра (в родительном падеже)
\chair{математической кибернетики и компьютерных наук}

% Тема работы
\title{Оценка сложности префикс"=функции, Z"=функции и алгоритма КМП}

% Курс
\course{2}

% Группа
\group{251}

% Факультет (в родительном падеже) (по умолчанию "факультета КНиИТ")
\department{факультета компьютерных наук и информационных технологий}

% Специальность/направление код - наименование
% \napravlenie{02.03.02 "--- Фундаментальная информатика и информационные технологии}
% \napravlenie{02.03.01 "--- Математическое обеспечение и администрирование информационных систем}
% \napravlenie{09.03.01 "--- Информатика и вычислительная техника}
\napravlenie{09.03.04 "--- Программная инженерия}
% \napravlenie{10.05.01 "--- Компьютерная безопасность}

% Для студентки. Для работы студента следующая команда не нужна.
\studenttitle{студентки}

% Фамилия, имя, отчество в родительном падеже
\author{Потапкиной Маргариты Андреевны}

% Заведующий кафедрой 
\chtitle{доцент, к.\,ф.-м.\,н.}
\chname{С.\,В.\,Миронов}

% Руководитель ДПП ПП для цифровой кафедры (перекрывает заведующего кафедры)
% \chpretitle{
%     заведующий кафедрой математических основ информатики и олимпиадного\\
%     программирования на базе МАОУ <<Ф"=Т лицей №1>>
% }
% \chtitle{г. Саратов, к.\,ф.-м.\,н., доцент}
% \chname{Кондратова\, Ю.\,Н.}

% Научный руководитель (для реферата преподаватель проверяющий работу)
\satitle{старший преподаватель} %должность, степень, звание
\saname{М.\,И.\,Сафрончик}

% Руководитель практики от организации (руководитель для цифровой кафедры)
\patitle{доцент, к.\,ф.-м.\,н.}
\paname{И.\,А.\,Батраева}

% Руководитель НИР
\nirtitle{доцент, к.\,п.\,н.} % степень, звание
\nirname{И.\,А.\,Батраева}

% Семестр (только для практики, для остальных типов работ не используется)
\term{2}

% Наименование практики (только для практики, для остальных типов работ не
% используется)
\practtype{учебная}

% Продолжительность практики (количество недель) (только для практики, для
% остальных типов работ не используется)
\duration{2}

% Даты начала и окончания практики (только для практики, для остальных типов
% работ не используется)
\practStart{01.07.2024}
\practFinish{13.01.2024}

% Год выполнения отчета
\date{2026}

\maketitle

% Включение нумерации рисунков, формул и таблиц по разделам (по умолчанию -
% нумерация сквозная) (допускается оба вида нумерации)
\secNumbering

\tableofcontents

% Раздел "Обозначения и сокращения". Может отсутствовать в работе
% \abbreviations
% \begin{description}
%     \item ... "--- ...
%     \item ... "--- ...
% \end{description}

% Раздел "Определения". Может отсутствовать в работе
% \definitions

% Раздел "Определения, обозначения и сокращения". Может отсутствовать в работе.
% Если присутствует, то заменяет собой разделы "Обозначения и сокращения" и
% "Определения"
% \defabbr

\intro
В данном отчёте производится анализ сложности и описание особенностей префикс"=функции, Z"=функции и алгоритма КМП.

\section{Анализ асимптотики}

\subsection{Префикс"=функция}
\small
\inputminted[breaklines]{cpp}{./prefix_function.cpp}
\normalsize

Основные части кода, влияющие на асимптотику алгоритма:

\begin{itemize}
	\item внешний цикл по элементам строки "--- $O(n)$;
	\item внутренний цикл \textit{while} "--- $O(n)$ суммарно по всем итерациям внешнего цикла.
\end{itemize}

Здесь $n$ "--- длина исходной строки.

\textbf{Итого асимптотику алгоритма иожно оценить как $O(n)$.}

\subsection{Z"=функция}
\small
\inputminted[breaklines]{cpp}{./z_function.cpp}
\normalsize

Основные части кода, влияющие на асимптотику алгоритма:

\begin{itemize}
	\item внешний цикл по элементам строки "--- $O(n)$;
	\item внутренний цикл \textit{while} "--- $O(n)$ суммарно по всем итерациям внешнего цикла.
\end{itemize}

Здесь $n$ "--- длина исходной строки.

\subsection{КМП}
\small
\inputminted[breaklines]{cpp}{./kmp.cpp}
\normalsize

Основные части кода, влияющие на асимптотику алгоритма:

\begin{itemize}
	\item построение префикс"=функции по новой строке, полученной конкатенацией шаблона, специального символа и строки, в которой осуществляем поиск "--- $O(n + m)$;
	\item анализ построенной префикс"=функции и поиск позиций совпадения "--- $O(n)$.
\end{itemize}

Здесь:
\begin{itemize}
	\item[$\bullet$] $n$ "--- длина строки, в которой осуществляем поиск;
	\item[$\bullet$] $m$ "--- длина шаблона, который ищем.
\end{itemize}

\section{Описание особенностей}

\subsection{Префикс"=функция}
Префикс"=функция находит для каждого префикса строки длину его наибольшего супрефикса "--- строки, которая одновременно является и префиксом, и суффиксом исходной. Наивный алгоритм её построения работает за $O(n^3)$. Оптимальный алгоритм, работающий за $O(n)$, основывается на том, что на каждом шаге префикс"=функция возрастает не более, чем на единицу (это происходит, только если текущий символ совпадает с символом, стоящим сразу после конца предыдущего супрефикса, и тем самым расширяет его). Значение переменной $j$, указывающей на длину рассматриваемого супрефикса, между итерациями внешнего цикла не меняется, а суммарно оно увеличивается не более чем на $n - 1$ (максимальная длина несобственного префикса всей строки, который может являться и её суффиксом). В цикле же \textit{while} оно только уменьшается (т.к. по свойству префикс"=функции $pr[j] < j + 1 \rightarrow pr[j - 1] < j$), суммарно также не более чем на $n - 1$ (т.к. $j$ не может стать меньше 0).

\subsection{Z"=функция}
Z"=функция находит для каждой позиции длину наибольшей подстроки, начинающейся с данной позиции и совпадающей c префиксом всей строки. Наивный алгоритм её построения работает за $O(n^2)$, оптимальный же алгоритм, имеющий асимптотику $O(n)$, опирается на идею Z"=блоков: мы храним левую и правую границу оптимального блока (с наиболее правой из всех правых границ). Теперь, если $i$ левее правой границы, значение Z"=функции для этой позиции мы уже знаем, а если оно правее "--- досчитаем это значение наивным образом. Поскольку каждая позиция участвует в наивном подсчёте Z"=функции только один раз (впоследствии она попадает в оптимальный диапазон и поэтому игнорируется), алгоритм работает за $O(n)$. 

\subsection{КМП}
Данный алгоритм используется для поиска подстроки в строке. Он применяется, когда необходимо найти все вхождения некоторого образца в тексте (для поиска большого количества образцов одновременно используется, например, алгоритм Ахо"=Корасик). Шаблон и текст конкатенируются, между ними записывается символ"=разделитель, которого нет в исходном алфавите: это нужно для того, чтобы не искать совпадения, большие длины шаблона. Теперь если значение префикс"=функции равно длине шаблона, мы нашли вхождение, заканчивающееся в данной позиции (возникло совпадение длины $m$ с началом, т.е. как раз был целиком найден шаблон).

% После введения — серии \section, \subsection и т.д.

% \conclusion

% Отобразить все источники. Даже те, на которые нет ссылок.
% \nocite{*}

\bibliographystyle{ugost2008}
\bibliography{thesis}

% Окончание основного документа и начало приложений Каждая последующая секция
% документа будет являться приложением
\appendix

\end{document}
